\documentclass[11pt]{article}

%%%%%%%%%%%%%%Include Packages%%%%%%%%%%%%%%%%%%%%%%%%%%
\usepackage{xcolor}
\usepackage{mathtools}
\usepackage[a4paper, total={6in, 8in}, margin=1.25in]{geometry}
\usepackage{amsmath}
\usepackage{amssymb}
\usepackage{paralist}
\usepackage{rsfso}
\usepackage{amsthm}
\usepackage{wasysym}
\usepackage[inline]{enumitem}   
\usepackage{hyperref}
\usepackage{tocloft}
\usepackage{titlesec}
\usepackage{colortbl}
\usepackage{wrapfig}
\usepackage{pdfpages}
\usepackage[american]{circuitikz}
%%%%%%%%%%%%%%%%%%%%%%%%%%%%%%%%%%%%%%%%%%%%%%%%%%%%%%%%



\begin{document}
\Large \noindent \textbf{Physics 5C Discussion (Worksheet 11)}\\
\normalsize
Department of Physics and Astronomy, UCLA\\

\noindent TA: Jinyan Miao (jnmiu@ucla.edu)\\
Office hours: Thursday 2:30 - 4:30 PM (PST) over Zoom \\
Zoom ID: 963 5916 9873 (passcode: hiJinyan)\\
Click for the \href{https://ucla.zoom.us/j/96359169873?pwd=05p7ssCPwTbl8KfxoBNCOuaIbt68jW.1}{\underline{Zoom link}} and \href{https://sites.google.com/umich.edu/jmiu/worksheets}{\underline{solutions to the worksheets}}.\\
\noindent\rule{8cm}{0.4pt}\\

\noindent \textbf{Problem 1}:
Jinyan is trying to determine the work function (binding energy) $\phi$ of a mysterious metallic material. He uses two different lasers to shine two different monochromatic lights on the metal. The first laser that he uses is a helium-neon laser, which produces a red light with wavelength $\lambda_{r}=633\,$nm. The second laser is a nitrogen laser, which produces an ultraviolet light with wavelength $\lambda_u = 337\,$nm. The helium-neon laser is very powerful, and it can deliver $1919\,$J of energy per second onto the surface of the mysterious metal. The nitrogen laser is much weaker, and it can deliver only around $19$\,J of energy per second onto the surface of the metal. \\

\noindent\textbf{Part (a)} Find the frequencies of the lights emitted by the two lasers.\\

\noindent\textbf{Part (b)} Find the energy of the photons emitted by the lasers, in both eV and J. \\

\noindent\textbf{Part (c)} If it is found that the electrons on the metal surface are ejected with a very tiny kinetic energy $\text{KE}\approx 0$\,eV when the nitrogen laser is used, what will be the maximal kinetic energy of those ejected electrons (if any) if the helium-neon laser is used?\\

\noindent\textbf{Part (d)} If it is found instead that the electrons on the metal surface are ejected with a very tiny kinetic energy $\text{KE}\approx 0$\,eV when the helium-neon laser is used, what will be the maximal kinetic energy of those ejected electrons (if any) if the nitrogen laser is used?\\

\noindent\textbf{Part (e)} Based on part (d), find the de Broglie wavelength of the electron that has its maximal kinetic energy, ejected by the nitrogen laser.

\newpage
\noindent \textbf{Problem 2}: Jinyan has been addicted to playing the game called Factorio with his friends, where they build automated factories. One late-game resource is Uranium: Uranium processing mostly produces the isotope $_{92}^{238}$U, which is essential in the production of nuclear fuel and “atomic bombs.” The $_{92}^{238}$U atom is unstable: It will undergo a atomic decay chain, releasing energy in this process (and this is why we say that $_{92}^{238}$U is radioactive), until it ends up at the stable isotope $^{206}_{82}$Pb. The decay chain is a combination of $\alpha$, $\beta$, and $\gamma$ decays. \\

\noindent\textbf{Part (a)} One of the processes in this nuclear decay chain is a $\gamma$ decay, 
\begin{align*}
^{214}_{83}\text{Bi}^* \to {}^{214}_{83}\text{Bi} + \gamma\,.
\end{align*}
The superscript $*$ is a notation that indicates $^{214}_{83}\text{Bi}^*$ has its nucleus in an excited state, and the nucleus ``jump" to the ground state (with the corresponding atom denoted as ${}^{214}_{83}\text{Bi}$) by emitting a $\gamma$-ray. The wavelength of such $\gamma$-ray is measured to be $\lambda = 7.05\cdot 10^{-31}\,\text{m}$. Find the energy difference between the $^{214}_{83}\text{Bi}^*$ atom and the ${}^{214}_{83}\text{Bi}$ atom in this process. \\

\noindent\textbf{Part (b)} Do a quick Google search, and explain why the $\gamma$-ray released in part (g) is highly dangerous to humans. (Hint: What is the frequency of the emitted $\gamma$-ray?) Since the $\gamma$ decay process in part (g) is part of the decay chain of $^{238}_{92}$U, your answer is the main reason why the decay of radioactive substances like $^{238}_{92}$U can ``kill" people. \\

\noindent\textbf{Part (c)} The half-life of $^{238}_{92}$U is around $4.5\cdot 10^9$ years. We have $1$ gram of pure $^{238}_{92}$U. How long does it take for at least $0.75$ grams of it to decay into other substances? \footnote{The decay of a radioactive substance is a random process: A standard nuclear reactor uses $20$ tonnes of Uranium a year, which is on the scale of $10^{28}$ Uranium atoms. Therefore, even though the half-life of Uranium seems to be very long, it is still highly likely to see some decay events happening at every single second given the large amount of atoms in the reactor.}\\

\noindent\textbf{Part (d)} By that time computed in part (c), how does the activity of Uranium sample compare to its initial activity at the time when we have a whole gram of $^{238}_{92}$U?



\newpage
\section*{Problem 1}
The key idea in this problem is that the energy of a single photon depends only on the light frequency (or wavelength), not on how powerful the laser is overall. A more powerful laser just sends more photons per second, but it does not make each individual photon more energetic. For the photoelectric effect, the important equation is
\begin{align*}
\text{KE}_{\max} = hf - \phi\,,
\end{align*}
where $h$ is Planck's constant and $\phi$ is the work function of the metal. We will use
\begin{align*}
h = 6.63\times 10^{-34}\,\text{J}\cdot \text{s}\,,\qquad
c=3.0\times 10^8\,\text{m/s}\,,\qquad
1\,\text{eV}=1.6\times 10^{-19}\,\text{J}\,.
\end{align*}

\noindent\textbf{Part (a)} The frequency of light is related to the wavelength by
\begin{align*}
f = \frac{c}{\lambda}\,.
\end{align*}
For the red helium-neon laser,
\begin{align*}
f_r = \frac{3.0\times 10^8\,\text{m/s}}{633\times 10^{-9}\,\text{m}}
= 4.74\times 10^{14}\,\text{Hz}\,.
\end{align*}
For the ultraviolet nitrogen laser,
\begin{align*}
f_u = \frac{3.0\times 10^8\,\text{m/s}}{337\times 10^{-9}\,\text{m}}
= 8.90\times 10^{14}\,\text{Hz}\,.
\end{align*}

\noindent\textbf{Part (b)} The energy of one photon is
\begin{align*}
E = hf\,.
\end{align*}
For the red photon,
\begin{align*}
E_r &= (6.63\times 10^{-34})(4.74\times 10^{14})\,\text{J}\\
&= 3.14\times 10^{-19}\,\text{J}\,.
\end{align*}
Converting this to eV,
\begin{align*}
E_r = \frac{3.14\times 10^{-19}\,\text{J}}{1.6\times 10^{-19}\,\text{J/eV}}
= 1.96\,\text{eV}\,.
\end{align*}

For the ultraviolet photon,
\begin{align*}
E_u &= (6.63\times 10^{-34})(8.90\times 10^{14})\,\text{J}\\
&= 5.90\times 10^{-19}\,\text{J}\,,
\end{align*}
so in eV,
\begin{align*}
E_u = \frac{5.90\times 10^{-19}\,\text{J}}{1.6\times 10^{-19}\,\text{J/eV}}
= 3.69\,\text{eV}\,.
\end{align*}
Therefore,
\begin{align*}
E_r = 1.96\,\text{eV} = 3.14\times 10^{-19}\,\text{J}\,,\qquad\qquad
E_u = 3.69\,\text{eV} = 5.90\times 10^{-19}\,\text{J}\,.
\end{align*}

\noindent\textbf{Part (c)} If the nitrogen laser ejects electrons with $\text{KE}\approx 0$\,eV, then the ultraviolet photons are just barely able to free the electrons. That means the work function of the metal must be
\begin{align*}
\phi \approx E_u = 3.69\,\text{eV}\,,
\end{align*}
because $\text{KE}_{\max} \approx 0 = E_u- \phi$. Now compare this with the energy of a red photon:
\begin{align*}
E_r = 1.96\,\text{eV} < \phi = 3.69\,\text{eV}\,.
\end{align*}
So a red photon does not even have enough energy to overcome the work function. Therefore \underline{no electrons are ejected at all} by the helium-neon laser, and hence there is no maximal kinetic energy to compute.\\

This is the important physics point: Even though the helium-neon laser is much more powerful overall, each red photon still has too little energy to eject an electron from this metal. Whether or not an electron can be ejected from the metal does not depend on how ``powerful" the laser is. It only depends on how ``powerful" each photon is. This is a straightforward comparison between the wave nature (laser power) and particle nature (photon energy) of light.\\

\noindent\textbf{Part (d)} Now suppose instead that the helium-neon laser ejects electrons with $\text{KE}\approx 0$\,eV. Then the red photons are just at the threshold, so the work function must be
\begin{align*}
\phi \approx E_r = 1.96\,\text{eV}\,,
\end{align*}
because $\text{KE}_{\max} \approx 0 =E_r - \phi$. For the nitrogen laser, the maximal kinetic energy is then
\begin{align*}
\text{KE}_{\max} = E_u - \phi = 3.69\,\text{eV} - 1.96\,\text{eV} = 1.73\,\text{eV}\,.
\end{align*}
Converting this to Joules,
\begin{align*}
\text{KE}_{\max} = (1.73)(1.6\times 10^{-19})\,\text{J}
= 2.77\times 10^{-19}\,\text{J}\,.
\end{align*}
So the maximal kinetic energy of the ejected electrons is
\begin{align*}
\text{KE}_{\max} = 1.73\,\text{eV} = 2.77\times 10^{-19}\,\text{J}\,.
\end{align*}

\noindent\textbf{Part (e)} The de Broglie wavelength is
\begin{align*}
\lambda = \frac{h}{mv}\,.
\end{align*}
To find the momentum $mv$ of the electron, we use the non-relativistic kinetic energy formula for electron with mass $m_e$, 
\begin{align*}
\text{KE} = \frac{1}{2}m_ev^2\,,
\end{align*}
so the momentum $m_ev$ of the electron is
\begin{align*}
m_ev = \sqrt{2m_e\text{KE}}\,.
\end{align*}
The kinetic energy is the one we found in part (c) and the mass is the electron mass,
\begin{align*}
\text{KE} = 2.77\times 10^{-19}\,\text{J}\,,\qquad
m_e = 9.11\times 10^{-31}\,\text{kg}\,.
\end{align*}
Therefore, we find the momentum of the electron to be
\begin{align*}
m_ev = 7.10\times 10^{-25}\,\text{kg}\cdot\text{m/s}\,.
\end{align*}
Now plug this into the de Broglie formula to get the de Broglie wavelength of electron,
\begin{align*}
\lambda = \frac{h}{m_ev}=\frac{6.63\times 10^{-34}}{7.10\times 10^{-25}}\,\text{m}
= 9.34\times 10^{-10}\,\text{m}= 0.93\,\text{nm}\,.
\end{align*}


\section*{Problem 2}

\noindent\textbf{Part (a)} Using the wavelength written in the problem, the emitted $\gamma$-photon has energy
\begin{align*}
E = hf = \frac{hc}{\lambda}\,.
\end{align*}
So we calculate
\begin{align*}
E &= \frac{(6.63\times 10^{-34}\,\text{J}\cdot \text{s})(3.0\times 10^8\,\text{m/s})}{7.05\times 10^{-31}\,\text{m}}= 2.82\times 10^5\,\text{J}\,.
\end{align*}
The $^{214}_{83}\text{Bi}^*$ atom decays to the $^{214}_{83}\text{Bi}$ atom by emitting the photon with energy calculated above. Energy has to be conserved, therefore,
\begin{align*}
\Delta E = 2.82\times 10^5\,\text{J}
\end{align*}
is the energy difference between the two states of the Bismuth atom. If we convert this into electron-Volts,
\begin{align*}
\Delta E = \frac{2.82\times 10^5\,\text{J}}{1.6\times 10^{-19}\,\text{J/eV}}
= 1.76\times 10^{24}\,\text{eV}\,.
\end{align*}

\noindent\textbf{Part (b)} First find the frequency of the emitted $\gamma$-ray:
\begin{align*}
f = \frac{c}{\lambda}
= \frac{3.0\times 10^8\,\text{m/s}}{7.05\times 10^{-31}\,\text{m}}
= 4.26\times 10^{38}\,\text{Hz}\,.
\end{align*}
This is an extremely high frequency, which means each photon has extremely high energy. An electromagnetic wave (or photon) with such a high frequency is usually referred as a $\gamma$-ray. The high-frequency nature of $\gamma$-rays is what makes it dangerous to humans. These $\gamma$-rays are called ``ionizing radiation." In other words, they can knock electrons out of atoms and molecules inside the human body, similar to the role of $\gamma$-rays in the photoelectric effect, where they knock electrons out of a metal by overcoming the binding energy. Therefore, these $\gamma$-rays can break chemical bonds, damage cells, and damage DNA. Also, $\gamma$-rays are very penetrating, so they can travel deep into the body instead of being stopped at the skin. Therefore, exposure to $\gamma$-rays can lead to severe tissue damage, cell death, mutations, and cancer. This is the main reason radioactive substances such as $^{238}_{92}\text{U}$ can be deadly.\\

\noindent\textbf{Part (c)} A half-life means the amount that remains is ``cut in half" after one half-life period of time. We start with $
1\,\text{g}$ of $^{238}_{92}\text{U}$. After one half-life, the amount remaining is half of $1\,$g, which is $0.50\,\text{g}\,.$ After two half-lives, the amount remaining is half of $0.50\,$g (that is, the count starts from the remaining part after one half-life), which is $
0.25\,\text{g}\,.$
That means the amount that has decayed into other types of substances is
\begin{align*}
1.00\,\text{g} - 0.25\,\text{g} = 0.75\,\text{g}
\end{align*}
after a period of two half-lives. So we need exactly two half-lives. Since one half-life is $4.5\times 10^9$ years, the required time is
\begin{align*}
\Delta t = 2\cdot (4.5\times 10^9\,\text{years}) = 9.0\times 10^9\,\text{years}\,.
\end{align*}
It takes $9.0\times 10^9\,\text{years}$ for at least $0.75\,$g of the original Uranium-238 to decay.\\

\noindent\textbf{Part (d)} Recall that, just like the number of radioactive atoms, the activity decays exponentially. If $R_0$ is the initial activity, $R$ is the activity after a time $t$ has passed, and $t_{1/2}$ is the half-life of Uranium-238, then
\begin{align*}
R=R_0\left(\frac{1}{2}\right)^{t/t_{1/2}}\,.
\end{align*}
From part (c), the time that has passed and the half-life of Uranium-238 are
\begin{align*}
t=9.0\times 10^9\,\text{years}\,,\qquad
t_{1/2}=4.5\times 10^9\,\text{years}\,.
\end{align*}
Plugging this result into the activity-decay formula, we find
\begin{align*}
R=R_0\left(\frac{1}{2}\right)^{t/t_{1/2}}=R_0\left(\frac{1}{2}\right)^2=\frac{R_0}{4}=0.25R_0\,.
\end{align*}
Thus, at the time computed in part (c), the activity of the Uranium sample is one-fourth of its initial activity, or $25\%$ of its initial activity. In other words, the activity has decreased by a factor of four. Just as expected, the amount of Uranium-238 reduces by a factor of 4, and the activity also reduces by a factor of 4. 




\end{document}
